\documentclass{grad-statement}

\addbibresource{refs.bib}

\usepackage{ali-todo}
\usepackage{xstring}

\newcommand\program{Cornell}

\author{Pattison}
\title{Statement of Purpose}

% CORNELL:
% Please use the Academic Statement of Purpose to describe (within 1000 words) the substantive research questions you are interested in pursuing during your graduate studies, and explain how our program would help you achieve your intellectual goals. Additionally, detail your academic background, intellectual interests, and any training or research experience you have received that you believe has prepared you for our program. Within your statement, please also identify specific faculty members whose research interests align with your own interests

\begin{document}

\noindent
I am applying to Cornell in pursuit of becoming a computer science professor conducting policy-relevant research at the intersection of economics, statistics, and computer science.

\section{Background and Prior Research}

I graduated \emph{summa cum laude} from Carleton College with a double major in mathematics and statistics.
Starting in the fall of my junior year, I worked on an applied cryptography project with Professor Nicholas Hopper at the University of Minnesota, developing a method for users of encrypted messaging apps to report abusive or misinformative content while still retaining privacy guarantees.
Previous research had struck a balance between accountability and privacy by proposing to attach a small token to every message that would allow a moderator to pinpoint a message's sender if the message was reported by its recipient.
I devised a way to democratize this process by requiring consensus among a panel of moderators before revealing the sender's identity and presented this construction as preliminary work at a major conference.

The following summer, I returned to Professor Hopper's group and strengthened my construction by having senders maintain anonymity until messages were reported many times or by many distinct users.
My solution used a novel cryptographic construction to accumulate reported messages without revealing anything until thresholds were exceeded, and used a zero-knowledge proof to convince recipients the attached metadata was valid.
I am preparing a manuscript on these results and hope to submit it for publication soon.

At Carleton, I also worked with Professor David Liben-Nowell on a paper studying a novel algorithmic problem inspired by clustering survey responses.
We show both that this problem is hard and that the hardness remains even in very restricted settings.
I individually contributed the main technical result: a proof that a naïve greedy algorithm achieves an approximation ratio of $0.5$ on a restricted class of problem instances.

I also served as a computational research assistant for an international physics research group and wrote two undergraduate senior theses.
For my math major, I presented an exposition of Dirichlet's unit theorem---a seminal result in algebraic number theory.
For my statistics major, I modeled sports games as sequences of discrete events to which I then applied deep-learning techniques originally developed for natural language processing.
I found that these position-aware sequence models provided little benefit over more classic statistical methods, suggesting that ``momentum'' does not play as large role in sports as common wisdom would suggest.

My senior year, I developed an interest in economics after taking an upper-level price-theory class where I was impressed by its ability to describe situations I had though were unnameable to rigorous quantitative study.
Since graduating, I've worked as a research assistant for Opportunity Insights---a group at Harvard that uses administrative microdata to study social mobility and its determinants in the United States.

At OI, I have been primarily involved with a project studying changes in American life expectancy over the last several decades and how and why those changes have been distributed geographically and across the income distribution.
We find that while overall mortality has significantly improved over the past two decades, nearly all of that improvement has been concentrated at the upper end of the income distribution.
Poor people, especially those in rural areas, have seen little improvement.
In some cases, there have even been declines.

One of my major contributions to this project has been the authorship of a pipeline to produce local estimates of life expectancy by income, geography, and other demographic characteristics using administrative microdata from the IRS, Census Bureau, and Social Security Administration.
While conceptually straightforward, this process has been plagued by statistical and computational challenges that arise when simultaneously estimating millions of parameters.
For example, how do you reliably estimate life expectancy in a population with only ten deaths?
How do you feasibly bootstrap standard errors for millions of estimates when a closed-form variance expression does not exist?
My solutions to these problems have made me realize that code---often seen by applied economists as an unavoidable but ultimately inconsequential nuisance---is in fact integral to the research process.
For example, my rewriting of part of our codebase improved processing speeds by orders of magnitude and shortened the research loop from several days to under an hour.

While not explicitly computer science, my work at OI has helped convince me that the questions that interest me most---for example opportunity, health inequality, and growing political polarization---can indeed by studied by rigorous quantitative methods.

\section{Future Research Interests}

At Cornell, I hope to continue the interests I've developed in college by pursuing research at the intersection of computer science, statistics, and economics.
This may take several forms.

The first is algorithms---both their theoretical analysis and their interactions with the real world.
Cornell has a very strong theory and algorithms group, and there are several faculty in particular with whom this sort of work could be done.
Both Allison Koenecke's and Jon Kleinberg's recent work on algorithmic fairness fall under this umbrella.

I'm also broadly interested in quantitative and computational social science and work that uses tools from computer science to study societal systems.
For example, nearly all of Nikhil Garg's and Jon Kleinberg's recent research is of interest to me.
In the IS department, I'm interested in working with Michael Macy to continue my ongoing research on growing polarization in the American public \footfullcite{pattison-2025-polarization}.

\medskip

More generally, I expect that my interests will to continue to evolve over the course of my PhD.
I'd benefit enormously from Cornell's excellence across many subfields of computer science and beyond in Information Science, ORIE, and Economics.

\end{document}
